\section{평균 Jacobian을 endpoint와 이차 나머지로 옮긴다}\label{sec:endpoint}
$y=g_2$, $J=D\HF(y)$로 두자. 원 convex $X$에서 $\HF\in C^2$이고 $y,g_F\in X$이므로 $y-sd$ 전체가 정의역에 놓인다. 기존의 정확한 평균식은
\begin{equation}
\overline A d=r,\qquad
\overline A=\int_0^1D\HF(y-sd)\dd s.
\label{eq:meanJ}
\end{equation}
평균 Jacobian의 값 자체를 새로 풀 필요 없이, 그 endpoint 차이를 적분형 Hessian으로 나타낼 수 있다.\cite{endpoint16}

\begin{theorem}[endpoint Taylor 전달]\label{thm:K}
다음 선형사상과 벡터를 정의한다.
\begin{equation}
K(d)v=\int_0^1(1-s)D^2\HF(y-sd)[d,v]\dd s,
\qquad R_2(d)=K(d)d.
\label{eq:K}
\end{equation}
그러면
\begin{equation}
\boxed{\overline A=J-K(d),\qquad Jd=r+R_2(d)}
\label{eq:JK}
\end{equation}
이며, 원 fixed 전처리 행렬 $C$에 대해
\begin{equation}
\boxed{d=Cr+Md+CR_2(d),\qquad M=I-CJ}
\label{eq:transfer}
\end{equation}
가 정확히 성립한다.
\end{theorem}
\begin{proof}
$D\HF(y-sd)=J-\int_0^sD^2\HF(y-td)[d,\cdot]\dd t$를 식 \eqref{eq:meanJ}에 대입한다. 적분 순서를 바꾸면 두 번째 적분의 가중치가 $1-t$가 되므로 $\overline A=J-K(d)$다. 따라서 $(J-K)d=r$이며, 왼쪽에 $C$를 곱하고 $CJ=I-M$을 사용한다. 같은 부호는 $0=\HF(y-d)=r-Jd+R_2$에서도 직접 확인된다.
\end{proof}

\subsection{HI opacity의 정확한 유리식 나머지}
활성 group마다
\begin{equation}
\ell_j=2ha_{2j},\quad D_j=1+\ell_j(1-x_2),\quad
D_{Fj}=D_j+\ell_jd_x=1+\ell_j(1-x_F)>0
\label{eq:DF}
\end{equation}
로 쓴다. 식 \eqref{eq:fullres}의 photon 항은 상수항을 제외하면 $\sum_j\nu_jQ_{Fj}/D_j$다. 따라서
\begin{align}
J_{ph}&=\sum_j\frac{\ell_jQ_{Fj}}{D_j^2}\nu_j\exx^{\mathsf T},&
\overline A_{ph}&=\sum_j\frac{\ell_jQ_{Fj}}{D_jD_{Fj}}\nu_j\exx^{\mathsf T},\\
K_{ph}(d)&=\left[\sum_j\frac{\ell_j^2Q_{Fj}d_x}{D_j^2D_{Fj}}\nu_j\right]\exx^{\mathsf T},&
R_{2,ph}(d)&=\sum_j\frac{\ell_j^2Q_{Fj}d_x^2}{D_j^2D_{Fj}}\nu_j.
\label{eq:rankone}
\end{align}
이는 모든 gas 성분을 묶는 일반 행렬의 추정이 아니라, 정확히 $x$ 한 열을 갖는 연산자다. 다음 단일 유리식 항등식이 식 \eqref{eq:rankone}을 증명한다.
\begin{equation}
\frac1{D+\ell d_x}-\frac1D+\frac{\ell d_x}{D^2}
=\frac{\ell^2d_x^2}{D^2(D+\ell d_x)}.
\label{eq:rationalid}
\end{equation}
물리적 $Q_{Fj}\ge0$에서 활성 group의 $\nu_{j,x},\nu_{j,w}\ge0$이므로 $R_{2,ph,x},R_{2,ph,w}\ge0$이다. 광자 guard와 inactive group에서는 이 항이 0이다. 다만 $K_{ph}$의 계수에는 $d_x$가 남아 있어, $R_{2,ph}$의 비음수성과 같은 부호 주장은 별도로 해야 한다.

기체 source의 나머지는
\begin{align}
R_{2,\BB}&=-h\widehat H_{\BB}[d,d],\qquad
K_{\BB}(d)v=-h\widehat H_{\BB}[d,v],\\
\widehat H_{\BB}&=2\int_0^1(1-s)D^2\BB_2(y-sd)\dd s
\in[D^2\BB_2](X),\notag\\
D^2\BB_2&=D^2F_2^{np}+\lambda\nu_H\otimes D^2q_2.
\label{eq:sourceHessian}
\end{align}
팽창은 $w$에 선형이므로 Hessian이 없다. 식 \eqref{eq:rankone}과 \eqref{eq:sourceHessian}은 source별 이차 나머지를 직접 구성하는 경로를 제공한다. 이번 수치 적용은 원 Full residual의 signed Hessian 전체를 사용한 보편 상계도 함께 계산하여 이 경로를 포함한다.

\section{원 componentwise majorant와 Hessian 상수}\label{sec:bounds}
\subsection{같은 \texorpdfstring{$C,B,q,S$}{C,B,q,S}를 유지한 첫 포함}
원 arithmetic 파일은 모든 $g\in X$와 고정 계열에서
\begin{equation}
|I-CD\HF(g)|\le B,\qquad
\nI{S^{-1}BS}\le q<1
\label{eq:contract}
\end{equation}
를 포함한다. 여기서
\begin{align}
S&=\diag(\ex(0.015),\ex(0.01),\ex(0.01),0.5),\notag\\
q&=\ex(0.00011736830563768177).
\label{eq:Sq}
\end{align}
$q$의 원 bits는 \texttt{3f1ec474231fbe16}이다. $\nS{v}:=\nI{S^{-1}v}$는 혼합 gas 단위를 제거한 노름이다. Neumann 급수에 의해 $CJ$와 $C\overline A$가 가역이고, 정사각 행렬 $C,J,\overline A$도 가역이다.\cite{phys14,Rump2010}

첫 포함은 scalar $q$만으로 반경을 키우지 않고 원 행렬 $B$를 사용한다. $u=\magv([Cr])$를 componentwise magnitude로 두고
\begin{equation}
(I-B)m=u,\qquad m\ge0
\label{eq:majorsolve}
\end{equation}
를 정확한 유리수 산술로 풀면 $|d|\le m$이다. 실제로 $d=Cr+(I-C\overline A)d$이므로 $|d|\le u+B|d|$이고, $(I-B)^{-1}=\sum_{k\ge0}B^k\ge0$다. 따라서
\begin{equation}
D_0=[Cr]+[-Bm,Bm]
\label{eq:D0}
\end{equation}
가 유효하다. 실제 실행은 원 signed $[I-C\overline A]$와의 교차를 추가했다. 이 첫 단계의 포함 반경은 이후 endpoint 보정의 입력이며, 새 물리 근 계산이 아니다.

\subsection{Scaled signed Hessian의 구성}
원 Full residual Hessian의 각 성분을 $[H^a_{jk}]=[D^2\HF]^a_{jk}(X)$라 하자. $S$의 대각성분 $s_i$로
\begin{equation}
\mathcal T_{ijk}=\frac{s_js_k}{s_i}\sum_aC_{ia}[H^a_{jk}],\qquad
L_H=\max_i\sum_{j,k}\magv(\mathcal T_{ijk})
\label{eq:LH}
\end{equation}
를 구성한다. $C$와 signed Hessian을 먼저 수축한 후 magnitude를 취하는 순서가 원 부호 정보를 유지한다. 그러면
\begin{equation}
\nI{S^{-1}CD^2\HF(g)[Su,Sv]}
\le L_H\nI{u}\nI{v},\qquad g\in X.
\label{eq:hbound}
\end{equation}
원 \path{SOURCE_BOX.json/G_gg}는 이미 $-2h$를 포함하는 \emph{Full residual} Hessian이다. 반면 \path{SOURCE_BOX_COMPONENTS/nonphoto}와 \path{SOURCE_BOX/q}는 source Hessian이므로 식 \eqref{eq:sourceHessian}에서 시간 계수를 적용해야 한다. 이 구별을 유지한 실제 읽기는
\begin{equation}
L_H\in[8.27802481433761309,8.27802481433761310]\times10^{-5}
\label{eq:LHvalue}
\end{equation}
를 얻었다.\cite{endpoint16}

\begin{proposition}[Endpoint 이차 오차]\label{prop:quad}
$\nS{d}\le a$이고 $\kappa=(1-q)^{-1}$라면
\begin{align}
\nI{S^{-1}CK(d)S}&\le\tfrac12L_Ha,&
\nS{CR_2(d)}&\le\tfrac12L_Ha^2,\\
\nS{d-J^{-1}r}&\le\frac{L_Ha^2}{2(1-q)}.
\label{eq:quadbound}
\end{align}
\end{proposition}
\begin{proof}
식 \eqref{eq:hbound}를 식 \eqref{eq:K}에 적분하면 $\int_0^1(1-s)\dd s=1/2$에서 첫 두 식을 얻는다. 식 \eqref{eq:JK}의 차이는 $d-J^{-1}r=(CJ)^{-1}CR_2$이며, 식 \eqref{eq:contract}에서 $\|(CJ)^{-1}\|_S\le\kappa$다.
\end{proof}

원 $X$ 전체의 Hessian을 유지한 이차 계수 $L_H/[2(1-q)]$는
\[
[4.13949825306495892,4.13949825306495893]\times10^{-5}
\]
에 포함된다.
이 계수는 실제 $d$의 작은 반경을 제곱해 사용한다. $X$ 자체의 크기를 제곱한 오차와 구별된다.

\section{유한 Neumann 보정과 닫힌 나머지}\label{sec:neumann}
\begin{theorem}[유한 endpoint 전달]\label{thm:neumann}
$P_m(M)=\sum_{k=0}^mM^k$, $v_m=P_m(M)Cr$로 두면
\begin{equation}
d=v_m+P_m(M)CR_2(d)+M^{m+1}d
\label{eq:finiteN}
\end{equation}
가 정확하며,
\begin{equation}
\boxed{\nS{d-v_m}\le\epsilon_m(a):=
q^{m+1}a+\frac12L_Ha^2\sum_{k=0}^mq^k.}
\label{eq:eps}
\end{equation}
\end{theorem}
\begin{proof}
식 \eqref{eq:transfer}에 $P_m$을 곱한다. $P_m(I-M)=I-M^{m+1}$이므로 식 \eqref{eq:finiteN}을 얻는다. $\|M\|_S\le q$와 명제 \ref{prop:quad}를 각 유한 항에 적용하면 식 \eqref{eq:eps}다. 무한 합의 수치 절단이나 새 Jacobian 역행렬의 근사값이 전제에 들어가지 않는다.
\end{proof}

실제 $J(y)$의 포함은 저장된 중심 Jacobian $J_c$와 같은 $X$의 Hessian을 이용한다.
\begin{equation}
J(y)=J_c+\int_0^1D^2\HF(c+t(y-c))[y-c,\cdot]\dd t.
\label{eq:Jaccenter}
\end{equation}
여기에는 $1-t$ 가중치가 없다. 전체 $X$의 Jacobian 포함과 식 \eqref{eq:Jaccenter}를 교차해 $[J(y)]$를 좁히고, 모든 sector에 같은 $[M]=I-C[J]$를 사용했다. 최종 포함은
\begin{equation}
D_{new}=D_0\cap\bigl([v_m]+S[-\epsilon_m,\epsilon_m]^4\bigr)
\label{eq:intersection}
\end{equation}
이다. 이번 실행은 $m=2$까지만 계산했다. 산술 오차와 이론 나머지를 모두 포함한 뒤 얻은 추가 반경이 실제 interval 폭보다 이미 매우 작기 때문이다.

\section{유한 관측량과 근사 adjoint}\label{sec:adjoint}
\subsection{선형 전자 수와 유리식 온도}
같은 마지막 시각에서 $\ell_e=(1,f_2,2f_2,0)$라 두면
\begin{equation}
\Delta e=\ell_ed,\quad \Delta n_e=n_2\ell_ed,\quad
\Delta w=\ew^{\mathsf T}d,\quad \Pi_F=\Pi_2-\ell_ed>0.
\label{eq:readout}
\end{equation}
온도는 유한 분모를 유지하면 정확히 선형 분자의 비가 된다.
\begin{equation}
n_T=C_T\ew^{\mathsf T}-T_2\ell_e,\quad N_T=n_Td,
\qquad \boxed{\Delta T=\frac{N_T}{\Pi_F}.}
\label{eq:finiteT}
\end{equation}
이는 $g_2$에서의 미분만 사용하는 온도 근사가 아니다. $T_F=C_T(w_2-d_w)/(\Pi_2-\ell_ed)$를 대입하여 얻은 항등식이다.

\begin{theorem}[근사 adjoint의 유한 보정]\label{thm:adjoint}
임의의 행벡터 $z^{\mathsf T}$와 dual residual $\eta^{\mathsf T}=n_T-z^{\mathsf T}J$에 대해
\begin{equation}
N_T=z^{\mathsf T}r+\eta^{\mathsf T}d+z^{\mathsf T}R_2(d).
\label{eq:dualresidual}
\end{equation}
특히 $z_m^{\mathsf T}=n_TP_m(M)C$라 두면 $\eta_m^{\mathsf T}=n_TM^{m+1}$이고,
\begin{equation}
|N_T-z_m^{\mathsf T}r|\le b_T\epsilon_m(a),\qquad
b_T=\|n_TS\|_1=T_2(s_x+f_2s_{y_1}+2f_2s_{y_2})+C_Ts_w.
\label{eq:goalerror}
\end{equation}
\end{theorem}
\begin{proof}
$n_T=z^{\mathsf T}J+\eta^{\mathsf T}$에 식 \eqref{eq:JK}를 적용하면 식 \eqref{eq:dualresidual}이다. $z_m^{\mathsf T}J=n_TP_m(M)(I-M)=n_T(I-M^{m+1})$에서 dual residual 식을 얻는다. 또한 $z_m^{\mathsf T}r=n_Tv_m$이므로 $|n_T(d-v_m)|\le\|n_TS\|_1\nS{d-v_m}$이다.
\end{proof}

분모 $\Pi_F$를 같은 구간으로 유지하는 관측량 포함에서는 식 \eqref{eq:goalerror}를 $\inf\Pi_F$로 나누면 된다. 예측 endpoint $v$ 자체의 분모를 사용하는 경우에도 닫힌 항등식이 있다. $e=d-v$, $\Pi_v=\Pi_2-\ell_ev$, $T_v=C_T(w_2-v_w)/\Pi_v$로 두면
\begin{equation}
\Delta T-\Delta T_v
=\frac{(C_T\ew^{\mathsf T}-T_v\ell_e)e}{\Pi_v-\ell_ee}.
\label{eq:tempredictor}
\end{equation}
따라서 예측 분모의 차이도 남은 $e$로 함께 제어할 수 있다. 실제 보고 수치는 식 \eqref{eq:finiteT}의 같은 $\Pi_F$를 사용한다.

\section{하나의 adjoint로 sector 기여와 부호를 보증한다}\label{sec:sector}
식 \eqref{eq:rg}의 두 gas 항과 식 \eqref{eq:seven}의 일곱 photon 항을 $r=\sum_s r_s$로 쓴다. 같은 평균 Jacobian으로
\begin{equation}
d_s=\overline A^{-1}r_s,\quad N_s=n_Td_s,\quad
\Delta T_s=N_s/\Pi_F
\label{eq:sectordef}
\end{equation}
를 정의하면 $d=\sum_sd_s$, $N_T=\sum_sN_s$, $\Delta T=\sum_s\Delta T_s$가 점별로 성립한다. 각 sector는 전체 endpoint와 같은 응답 연산자 및 온도 분모를 쓴다.

\begin{theorem}[공통 sector 오차와 부호 여유]\label{thm:sector}
$c_s\ge\nS{Cr_s}$, $\kappa=(1-q)^{-1}$라 하자. $\eta_m$과 $z_m$를 정리 \ref{thm:adjoint}와 같이 두면
\begin{equation}
N_s-z_m^{\mathsf T}r_s=(\eta_m^{\mathsf T}+z_m^{\mathsf T}K)d_s
\label{eq:sectoridentity}
\end{equation}
이며,
\begin{equation}
|N_s-z_m^{\mathsf T}r_s|\le E_s,
\qquad E_s=\kappa c_sb_T\left(q^{m+1}
+\frac12L_Ha\sum_{k=0}^mq^k\right).
\label{eq:sectorerror}
\end{equation}
$[z_m^{\mathsf T}r_s]=[L_s,U_s]$라면 $L_s>E_s$는 양의, $U_s<-E_s$는 음의 sector 부호를 보증한다.
\end{theorem}
\begin{proof}
$(J-K)d_s=r_s$이므로 $Jd_s=r_s+Kd_s$다. 여기에 $n_T=z_m^{\mathsf T}J+\eta_m^{\mathsf T}$를 적용하여 식 \eqref{eq:sectoridentity}를 얻는다. 식 \eqref{eq:contract}를 평균한 결과에서 $\nS{d_s}\le\kappa c_s$다. $\|\eta_m^{\mathsf T}S\|_1\le b_Tq^{m+1}$이며,
\[
\|z_m^{\mathsf T}KS\|_1
=\|n_TP_m(M)CKS\|_1
\le b_T\sum_{k=0}^mq^k\frac12L_Ha.
\]
두 항에 $\nS{d_s}$를 곱하면 식 \eqref{eq:sectorerror}다. 마지막 부호 판정은 $\Pi_F>0$에서 분자와 온도 기여의 부호가 같다는 사실을 사용한다.
\end{proof}

부호가 구조적으로 상쇄되는 birth 항은 공통 평균 adjoint
$\overline z^{\mathsf T}=n_T\overline A^{-1}$의 componentwise 포함을 직접 사용하는 편이 유리하다. $c_C=\|S^{-1}CS\|_\infty$라 두면
\begin{equation}
\| (\overline z-z_m)^{\mathsf T}S\|_1\le\Theta_z,
\quad\Theta_z=\kappa c_Cb_T\left(q^{m+1}
+\frac12L_Ha\sum_{k=0}^mq^k\right).
\label{eq:adjointrow}
\end{equation}
실제로 $(\overline z-z_m)^{\mathsf T}\overline A=\eta_m^{\mathsf T}+z_m^{\mathsf T}K$이며 $\|\overline A^{-1}\|_S\le\kappa c_C$이므로 위 식을 얻는다. 각 성분에 $\Theta_z/s_i$를 더하고 빼면 $[\overline z_i]$를 포함한다. 이 공통 행벡터가 다음 절의 birth 부호를 닫는다.
